\documentclass[10pt,a4paper]{article}
\usepackage[utf8]{inputenc}
\usepackage[T1]{fontenc}
\usepackage[czech]{babel}
\usepackage{lmodern}
\usepackage{geometry}
\usepackage{xcolor}
\usepackage{enumitem}
\usepackage{hyperref}
\usepackage{titlesec}

\geometry{margin=1.55cm}
\hypersetup{
  colorlinks=true,
  linkcolor=blue!55!black,
  urlcolor=blue!55!black,
  pdftitle={BSS mini tahak ke statnicim},
  pdfauthor={Petr Marinec}
}
\setlength{\parindent}{0pt}
\setlength{\parskip}{2pt}
\setlist[itemize]{leftmargin=*, topsep=1pt, itemsep=1pt, parsep=0pt}
\titleformat{\section}{\large\bfseries}{\thesection}{0.5em}{}
\titleformat{\subsection}{\normalsize\bfseries}{\thesubsection}{0.5em}{}

\newcommand{\otazka}[2]{\section*{#1}\addcontentsline{toc}{section}{#1}\textit{#2}\par\medskip}
\newcommand{\must}[1]{\colorbox{blue!7}{\parbox{\dimexpr\linewidth-2\fboxsep}{\textbf{Minimum k odpovědi:} #1}}\par\medskip}
\newcommand{\blok}[1]{\subsection*{#1}}

\title{BSS mini tahák ke státnicím\\\large otázky 15--23}
\author{248689}
\date{Aktualizováno 18. června 2026}

\begin{document}
\maketitle
\tableofcontents
\newpage

\section*{Jak to používat}
\addcontentsline{toc}{section}{Jak to používat}
Toto není učebnice, ale rychlá osnova odpovědi. U každé otázky řekni nejdřív hlavní vymezení, potom projdi přesně podotázky ze zadání a zakonči jedním praktickým příkladem.

\newpage

\otazka{15. Vymezení bezpečnostních studií}{Bezpečnostní studia, bezpečnost, hrozba, riziko, vnitřní/vnější bezpečnost, sektory, bezpečnostní politika a její analýza v kyberbezpečnosti.}

\blok{Vymezení bezpečnostních studií}
\begin{itemize}
  \item Zkoumají, co je ohroženo, kdo nebo co ohrožuje, jak hrozbu zvládat a jak obnovit bezpečnost.
  \item Původně hlavně válka a vojenská bezpečnost; dnes i vnitřní, ekonomická, společenská, environmentální, kriminální a kybernetická bezpečnost.
  \item Referenční objekt = co chráním: stát, společnost, infrastruktura, instituce, firma, člověk, data.
\end{itemize}

\blok{Bezpečnost, hrozba, zranitelnost, riziko}
\begin{itemize}
  \item Bezpečnost = hrozby jsou omezené na přijatelnou úroveň a systém je schopný je zvládat.
  \item Kyberbezpečnost = ochrana sítí, systémů, služeb, dat a uživatelů; hlavně důvěrnost, integrita, dostupnost, odolnost a obnova.
  \item Hrozba = možný zdroj škody; zranitelnost = slabina; riziko = pravděpodobnost a dopad využití hrozby; incident = událost, která už nastala.
  \item Příklad: ransomware je hrozba, nezálohovaný server je zranitelnost, výpadek nemocnice je dopad rizika.
\end{itemize}

\blok{Vnitřní a vnější bezpečnost}
\begin{itemize}
  \item Vnitřní = pořádek, kriminalita, krizové řízení, ochrana obyvatelstva a fungování institucí uvnitř státu.
  \item Vnější = obrana proti zahraniční agresi, nátlaku, špionáži a hybridnímu působení.
  \item V kyberprostoru se prolínají: útok ze zahraničí může dopadnout na českou nemocnici a řeší ho NÚKIB, policie, zpravodajské služby i diplomacie.
\end{itemize}

\blok{Sektory bezpečnosti}
\begin{itemize}
  \item Vojenský: obrana, ozbrojený konflikt, kyberobrana.
  \item Politický: instituce, suverenita, volby, důvěra ve stát.
  \item Ekonomický: finance, dodavatelské řetězce, strategické technologie.
  \item Societální: identita, soudržnost, dezinformace.
  \item Environmentální: přírodní hrozby, klima; u KI i all-hazards odolnost.
\end{itemize}

\blok{Bezpečnostní politika a analýza}
\begin{itemize}
  \item Bezpečnostní politika = rámec, kterým stát určuje, co chrání, před čím, kdo za to odpovídá a jakými nástroji.
  \item Analýza politiky: referenční objekt, hrozby, aktéři, nástroje, náklady, vedlejší dopady, měření účinnosti.
  \item V kyberbezpečnosti se ptám: jaká služba/data jsou chráněna, kdo útočí, jaké jsou zranitelnosti, jaké povinnosti a instituce reagují.
\end{itemize}

\must{Bezpečnostní studia = co chráním, před čím a jak. Hrozba není riziko. U kyberbezpečnosti vždy propojit techniku, instituce a právo.}

\newpage

\otazka{16. Bezpečnostní strategie}{Dokumenty ČR a význam kyberbezpečnosti v nich. Bezpečnostní systém ČR a role institucí v kyberbezpečnosti.}

\blok{Bezpečnostní strategie}
\begin{itemize}
  \item Strategie = politicko-strategický dokument; neurčuje jednotlivci povinnost jako zákon, ale říká směr, priority a cíle.
  \item Řetězec: analýza hrozeb -- strategie -- akční plán -- zákon/vyhláška/metodika -- praxe -- vyhodnocení.
  \item Příklad: strategie řekne chránit zdravotnictví proti ransomwaru; zákon a metodiky nastaví povinnosti, hlášení, zálohy a reakci.
\end{itemize}

\blok{Dokumenty ČR}
\begin{itemize}
  \item Bezpečnostní strategie ČR 2023: nejširší rámec; zhoršené prostředí, Rusko, hybridní a kybernetické hrozby, odolnost státu a společnosti.
  \item Obranná strategie ČR 2023: obrana státu, NATO, odstrašení, ozbrojené síly, vojenské působení v kybernetickém a informačním prostoru.
  \item Národní strategie kybernetické bezpečnosti 2026--2030: bezpečná strategická infrastruktura, celospolečenská připravenost, mezinárodní spolupráce.
  \item Strategie pro odolnost subjektů kritické infrastruktury 2026--2029: all-hazards odolnost základních služeb.
\end{itemize}

\blok{Bezpečnostní systém ČR}
\begin{itemize}
  \item Soubor institucí, pravidel a procesů pro prevenci, reakci a obnovu.
  \item Prezident: hlava státu a vrchní velitel ozbrojených sil, ne běžný řídicí orgán civilní kyberbezpečnosti.
  \item Parlament: zákony a kontrola vlády.
  \item Vláda: hlavní výkonný orgán bezpečnostní politiky.
  \item Bezpečnostní rada státu: pracovní orgán vlády pro koordinaci.
  \item Ústřední krizový štáb, ministerstva, kraje, obce a IZS: krizové řízení a sektorové úkoly.
\end{itemize}

\blok{Role institucí v kyberbezpečnosti}
\begin{itemize}
  \item NÚKIB: hlavní civilní autorita, regulace ZKB, GovCERT.CZ, varování, výstrahy, reaktivní protiopatření.
  \item CSIRT.CZ: národní CSIRT provozovaný CZ.NIC, praktická koordinace incidentů mimo státní GovCERT.
  \item Policie a státní zastupitelství: kyberkriminalita a trestní řízení.
  \item Zpravodajské služby: BIS, ÚZSI, VZ; špionáž, hrozby, zahraniční aktéři.
  \item MO, Vojenské zpravodajství, Armáda ČR: kyberobrana, vojenský rozměr.
  \item MZV: diplomacie, atribuce, sankce, EU/NATO reakce.
  \item DIA a ČTÚ: digitální služby státu, eGovernment, komunikace a infrastruktura.
  \item Soukromý sektor: vlastní a provozuje velkou část sítí a kritických služeb.
\end{itemize}

\must{Strategie určí směr, zákon povinnosti. U institucí rozliš civilní kyberbezpečnost, kyberkriminalitu, zpravodajství a vojenskou kyberobranu.}

\newpage

\otazka{17. Kyberválka}{Definice, historie, trendy, aktéři, atribuce, odstrašení, netwars v kyberkonfliktu.}

\blok{Definice}
\begin{itemize}
  \item Kyberválka = použití kyberoperací státem nebo politickým aktérem jako součást ozbrojeného konfliktu nebo s účinky srovnatelnými se závažným vojenským útokem.
  \item Ne každý státní kyberútok je kyberválka; většina útoků je špionáž, kriminalita, sabotáž, nátlak nebo podprahové působení.
  \item Rozliš: kyberbezpečnost, kyberobrana, kyberkriminalita, kybernetický konflikt, kyberválka.
\end{itemize}

\blok{Historie a příklady}
\begin{itemize}
  \item Morrisův červ 1988: raná ukázka zranitelnosti propojených sítí.
  \item Estonsko 2007: DDoS a tlak na stát, často uváděný příklad politického kyberkonfliktu.
  \item Gruzie 2008: kyberoperace souběžně s konvenční válkou.
  \item Stuxnet: sabotáž fyzického procesu přes kyberprostředky.
  \item Ukrajinská energetika 2015/2016: výpadky elektřiny.
  \item NotPetya 2017: destruktivní malware s globálními ekonomickými dopady.
  \item SolarWinds: dodavatelský řetězec a špionáž, ne automaticky kyberválka.
  \item Ukrajina po roce 2022: propojení kybernetických, informačních a konvenčních operací.
\end{itemize}

\blok{Současné trendy}
\begin{itemize}
  \item Podprahové operace: pod hranicí ozbrojeného útoku.
  \item APT skupiny: dlouhodobé, dobře financované a často státem podporované operace.
  \item Pre-positioning v kritické infrastruktuře: příprava přístupu pro budoucí krizi.
  \item Útoky na cloud, identity, telekomunikace a dodavatelský řetězec.
  \item AI zrychluje phishing, analýzu zranitelností a informační operace.
  \item Hacktivismus a proxy aktéři maskují vazbu na stát.
\end{itemize}

\blok{Aktéři}
\begin{itemize}
  \item Státy, armády a zpravodajské služby.
  \item Státem podporované APT skupiny.
  \item Proxy aktéři a hacktivisté.
  \item Kyberkriminální skupiny, které mohou být tolerované nebo využité státem.
  \item Soukromé firmy, dodavatelé, dobrovolníci.
\end{itemize}

\blok{Atribuce a odstrašení}
\begin{itemize}
  \item Atribuce = přisouzení útoku aktérovi; má technickou, operační, zpravodajskou a politicko-právní rovinu.
  \item IP adresa nestačí: proxy, VPN, botnety, false flag, sdílený malware, kompromitovaná infrastruktura.
  \item Odstrašení trestem: sankce, diplomacie, trestní stíhání, protiopatření.
  \item Odstrašení odepřením: útok se nevyplatí, protože systém je dobře chráněn.
  \item Odolnost: rychlá obnova snižuje efekt útoku.
\end{itemize}

\blok{Netwars}
\begin{itemize}
  \item Netwar = konflikt vedený síťově organizovanými aktéry; nejde jen o technické hackování.
  \item Typické: decentralizace, informační kampaně, koordinované DDoS, swarming.
  \item V kyberkonfliktu doplňuje kyberútoky propagandou, úniky dat, mobilizací online komunit a tlakem na veřejnost.
\end{itemize}

\must{Kyberválka je úzký pojem. Řekni definici, příklady vývoje, aktéry, proč je těžká atribuce, tři logiky odstrašení a netwar jako síťový konflikt.}

\newpage

\otazka{18. Ochrana kritické infrastruktury}{Kritická infrastruktura, kyberútoky na KI, přisouzení a odstrašování v kyberprostoru.}

\blok{Kritická infrastruktura}
\begin{itemize}
  \item Jde o služby a infrastrukturu, jejichž narušení ohrozí základní funkce státu, ekonomiku, bezpečnost nebo životy.
  \item Důraz není jen na data, ale na kontinuitu základní služby: elektřina, voda, zdravotnictví, doprava, finance, komunikace.
  \item CER/ZKI = all-hazards odolnost kritických subjektů; ZKB/NIS2 = kybernetická bezpečnost regulovaných služeb.
\end{itemize}

\blok{Kyberútoky na KI}
\begin{itemize}
  \item Ransomware: šifrování systémů, výpadek provozu, vydírání.
  \item Sabotáž OT/ICS: zásah do průmyslových systémů, fyzické dopady.
  \item DDoS: nedostupnost veřejných služeb.
  \item Supply chain: napadení dodavatele nebo aktualizace.
  \item Špionáž a pre-positioning: dlouhodobý přístup pro budoucí krizi.
\end{itemize}

\blok{Ochrana}
\begin{itemize}
  \item Řízení rizik, segmentace IT/OT, zálohy, MFA, monitoring, reakční plány.
  \item Cvičení, kontinuita provozu, krizová komunikace.
  \item Bezpečnost dodavatelů a smluvní požadavky.
  \item Hlášení incidentů a spolupráce s NÚKIB/GovCERT/CSIRT.
\end{itemize}

\blok{Atribuce a odstrašování}
\begin{itemize}
  \item U KI je atribuce politicky citlivá, protože může vést k diplomacii, sankcím nebo bezpečnostním opatřením.
  \item Odstrašení trestem: sankce, veřejná atribuce, trestní stíhání.
  \item Odstrašení odepřením: snížení šance útoku uspět.
  \item Odolnost: útok nezpůsobí strategický efekt, protože služba se rychle obnoví.
\end{itemize}

\must{U KI vždy řekni: nejde jen o kyberincident, ale o dopad na základní službu a společnost. Rozliš CER/ZKI a ZKB/NIS2.}

\newpage

\otazka{19. Právní úprava kyberbezpečnosti v ČR a EU}{Základní instituty, principy, povinné orgány, systém zajištění kybernetické bezpečnosti.}

\blok{EU rámec}
\begin{itemize}
  \item NIS2: základní směrnice pro kyberbezpečnost důležitých a základních subjektů; v ČR přes nový ZKB.
  \item DORA: finanční sektor, ICT rizika a digitální provozní odolnost.
  \item CRA: produkty s digitálními prvky, bezpečný vývoj a zranitelnosti.
  \item Cyber Solidarity Act: evropský systém kybernetického varování, cyber hubs, mechanismus kybernetické nouze, Evropská kybernetická rezerva, přezkum incidentů.
\end{itemize}

\blok{Český rámec}
\begin{itemize}
  \item Zákon o kybernetické bezpečnosti z roku 2025 transponuje NIS2 a nahrazuje starší úpravu.
  \item Regulovaná služba, režim vyšších/nižších povinností, sebeidentifikace subjektu.
  \item Bezpečnostní opatření: organizační a technická opatření podle rizika.
  \item Hlášení incidentů: včasné varování, oznámení, závěrečná zpráva podle režimu a vyhlášek.
\end{itemize}

\blok{Principy}
\begin{itemize}
  \item Risk-based approach: povinnosti se odvíjejí od rizika.
  \item Proporcionalita: opatření mají odpovídat významu služby a dopadu.
  \item Technologická neutralita: právo nemá být svázané s jednou technologií.
  \item Odpovědnost vedení: kyberbezpečnost není jen problém IT.
  \item Autonomie regulovaného subjektu: stát stanoví rámec, subjekt řídí vlastní rizika.
\end{itemize}

\blok{NÚKIB a nástroje}
\begin{itemize}
  \item NÚKIB: regulátor, metodika, dohled, GovCERT.CZ, evidence a reakce.
  \item Protiopatření podle ZKB: výstraha, varování, reaktivní protiopatření.
  \item Výstraha upozorňuje na existující nebo hrozící nebezpečí.
  \item Varování upozorňuje na hrozbu, kterou mají dotčení zohlednit v řízení rizik.
  \item Reaktivní protiopatření ukládá konkrétní technická nebo organizační opatření.
  \item Nápravná opatření při dohledu jsou něco jiného než protiopatření.
\end{itemize}

\blok{Povinné orgány a systém}
\begin{itemize}
  \item Manažer kybernetické bezpečnosti, architekt, výbor, garant aktiva, auditor nebo další role podle režimu.
  \item Systém řízení bezpečnosti informací, řízení aktiv, rizik, dodavatelů, incidentů, kontinuity a obnovy.
  \item Stav kybernetického nebezpečí vyhlašuje NÚKIB na úřední desce; nejvýše 30 dnů, celkem maximálně 60 dnů při prodloužení.
\end{itemize}

\must{NIS2 není sama český zákon. V ČR mluv hlavně o ZKB 2025, regulovaných službách, režimech povinností, NÚKIB, bezpečnostních opatřeních a hlášení incidentů.}

\newpage

\otazka{20. Kyberkriminalita}{Prameny práva, typická trestná činnost, klasifikace trestných činů, právní kvalifikace, postupy, kritéria a mezinárodní spolupráce.}

\blok{Prameny práva}
\begin{itemize}
  \item Národní: trestní zákoník, trestní řád, zákon o Policii ČR, zákon o elektronických komunikacích.
  \item EU: směrnice o útocích na informační systémy, e-evidence, evropský vyšetřovací příkaz, směrnice k platebním podvodům.
  \item Mezinárodní: Budapešťská úmluva, budoucí protokoly a nástroje OSN podle ratifikace.
\end{itemize}

\blok{Klasifikace}
\begin{itemize}
  \item Cyber-dependent: trestný čin existuje díky ICT, např. neoprávněný přístup, malware, DDoS.
  \item Cyber-enabled: tradiční kriminalita posílená ICT, např. podvod, vydírání, stalking.
  \item Cyber-supported: ICT slouží jako komunikace, úložiště nebo prostředek organizace.
\end{itemize}

\blok{Typická trestná činnost}
\begin{itemize}
  \item Neoprávněný přístup, poškození/zásah do dat a systému.
  \item Opatření a přechovávání přístupového zařízení nebo hesla.
  \item Podvod, phishing, platební podvody, krádež identity.
  \item Vydírání, ransomware, doxing, stalking, dětská pornografie.
  \item Porušení autorského práva, legalizace výnosů z trestné činnosti.
\end{itemize}

\blok{Právní kvalifikace}
\begin{itemize}
  \item Neříkat jen "hacking"; popsat skutek, chráněný zájem, úmysl, škodu, následek a kvalifikovanou skutkovou podstatu.
  \item U neoprávněného přístupu řešit, zda bylo překonáno bezpečnostní opatření.
  \item U ransomware kombinace: zásah do systému/dat, vydírání, škoda, případně KI nebo zdravotnictví jako přitěžující kontext.
\end{itemize}

\blok{Související postupy}
\begin{itemize}
  \item Zajištění dat, domovní prohlídka, vydání/odnětí věci, uchování dat.
  \item Metadata vs obsah komunikace: různé procesní tituly a různá intenzita zásahu.
  \item Forenzní postup: kopie, hash, protokol, chain of custody.
  \item Spolupráce s poskytovateli služeb a bankami.
\end{itemize}

\blok{Mezinárodní spolupráce}
\begin{itemize}
  \item MLA: klasická právní pomoc stát--stát.
  \item Evropský vyšetřovací příkaz: EU nástroj pro vyšetřovací úkony.
  \item EPOC: evropský příkaz k vydání uložených elektronických důkazů.
  \item EPOC-PR: evropský příkaz k uchování dat.
  \item 24/7 kontaktní síť podle Budapešťské úmluvy.
\end{itemize}

\must{Postupuj od faktů ke skutkové podstatě. Vždy řekni prameny, typ útoku, kvalifikaci, procesní zajištění důkazů a přeshraniční spolupráci.}

\newpage

\otazka{21. Elektronické důkazy a jejich zajišťování}{Procesní instituty, praktické využití, nakládání s elektronickými důkazy, elektronické dokumenty.}

\blok{Elektronický důkaz}
\begin{itemize}
  \item Informace v digitální podobě použitelná v řízení: soubor, log, e-mail, metadata, chat, cloud, telefon, záznam, blockchain stopa.
  \item Není automaticky slabší než listina; musí být prokázán původ, integrita, čas a zákonnost získání.
  \item Důkazní hodnota stojí na forenzní integritě a dokumentaci.
\end{itemize}

\blok{Procesní instituty}
\begin{itemize}
  \item Vydání věci/dat a odnětí věci/dat.
  \item Domovní prohlídka a prohlídka jiných prostor.
  \item Uchování dat podle trestního řádu.
  \item Odposlech a záznam telekomunikačního provozu = obsah budoucí komunikace.
  \item Zjištění údajů o telekomunikačním provozu = metadata.
  \item Operativně pátrací prostředky a součinnost poskytovatelů podle zákonného titulu.
\end{itemize}

\blok{Praktické využití}
\begin{itemize}
  \item Ransomware: zajistit zařízení, logy, výkupné, komunikaci, kryptopeněženky.
  \item Phishing: e-mailové hlavičky, URL, domény, bankovní transakce, přístupové logy.
  \item Cloud: rychle uchovat data, protože mohou být změněna nebo smazána.
  \item Mobil: forenzní image, hesla/biometrie podle zákonných limitů, dokumentace úkonů.
\end{itemize}

\blok{Nakládání s důkazy}
\begin{itemize}
  \item Pracovat s kopií, originál chránit.
  \item Hash pro integritu.
  \item Chain of custody: kdo, kdy, kde, co převzal a jak s tím nakládal.
  \item Oddělit obsah, metadata, provozní data a údaje o účastníkovi.
  \item Pozor na nezákonné získání, změnu dat, chybějící protokol nebo přerušený řetězec držení.
\end{itemize}

\blok{Přeshraniční důkazy}
\begin{itemize}
  \item EVP: orgán žádá orgán v EU.
  \item MLA: klasická právní pomoc.
  \item Budapešťská úmluva: uchování dat a 24/7 kontaktní síť.
  \item EPOC/EPOC-PR: přímější EU režim pro vydání nebo uchování uložených elektronických důkazů; není živý odposlech a neprolamuje šifrování.
\end{itemize}

\blok{Elektronické dokumenty}
\begin{itemize}
  \item Elektronický dokument může být důkaz stejně jako listina.
  \item Podpis/pečeť/časové razítko řeší původ, integritu a čas, ne pravdivost obsahu.
  \item U dokumentu řešit formát, metadata, podpis, certifikát, časové razítko a historii změn.
\end{itemize}

\must{Klíčové rozlišení: obsah budoucí komunikace, metadata, uchování existujících dat, vydání/odnětí věci. Pak říct hash, kopie, protokol, chain of custody.}

\newpage

\otazka{22. Ochrana osobních údajů}{Právní úprava, principy, zásady, rizika, proporcionalita, účely, právní tituly, ÚOOÚ.}

\blok{Právní úprava}
\begin{itemize}
  \item GDPR = hlavní obecný rámec.
  \item Zákon o zpracování osobních údajů = česká adaptační úprava a ÚOOÚ.
  \item Zvláštní předpisy: ePrivacy, pracovní právo, zdravotnictví, policie, kyberbezpečnost podle kontextu.
\end{itemize}

\blok{Základní pojmy}
\begin{itemize}
  \item Osobní údaj = informace o identifikované nebo identifikovatelné osobě.
  \item Zpracování = jakákoli operace s údaji: sběr, uložení, analýza, předání, výmaz.
  \item Správce určuje účely a prostředky; zpracovatel jedná pro správce.
  \item Pseudonymizace není anonymizace; pseudonymizovaná data zůstávají osobními údaji.
\end{itemize}

\blok{Zvláštní kategorie}
\begin{itemize}
  \item Citlivé údaje: zdraví, biometrie pro identifikaci, genetika, politické názory, náboženství, odbory, sexuální orientace.
  \item Nestačí jen právní titul podle čl. 6 GDPR; je potřeba i výjimka podle čl. 9 GDPR.
  \item Prakticky: zdravotnictví, biometrická docházka, profilování citlivých kategorií.
\end{itemize}

\blok{Zásady}
\begin{itemize}
  \item Zákonnost, korektnost, transparentnost.
  \item Účelové omezení: data jen pro jasný účel.
  \item Minimalizace: jen nezbytné údaje.
  \item Přesnost, omezení uložení, integrita a důvěrnost.
  \item Odpovědnost správce: musí doložit, že pravidla dodržel.
\end{itemize}

\blok{Právní tituly}
\begin{itemize}
  \item Souhlas: svobodný, konkrétní, informovaný, odvolatelný.
  \item Smlouva: zpracování nutné pro plnění smlouvy.
  \item Právní povinnost: účetnictví, zákonné evidence.
  \item Životně důležitý zájem: nouzové situace.
  \item Veřejný úkol / veřejná moc.
  \item Oprávněný zájem: nutný balancing test.
\end{itemize}

\blok{Oprávněný zájem}
\begin{itemize}
  \item Test účelu: jaký legitimní zájem sleduji.
  \item Test nezbytnosti: nejde to méně invazivně?
  \item Balancing test: nepřevažují práva a svobody subjektu údajů?
  \item Příklad: kamera u vstupu do skladu může obstát; kamera se zvukem na všechny zaměstnance kvůli výkonu spíš ne.
\end{itemize}

\blok{Riziko, DPIA, proporcionalita}
\begin{itemize}
  \item Riziko = dopad na práva a svobody lidí, ne jen technické riziko firmy.
  \item DPIA se dělá před rizikovým zpracováním, např. rozsáhlé profilování, citlivé údaje, systematické monitorování.
  \item Proporcionalita: vhodnost, potřebnost, přiměřenost v užším smyslu.
  \item Privacy by design/default: ochrana soukromí už v návrhu systému.
\end{itemize}

\blok{DPO, práva subjektů a ÚOOÚ}
\begin{itemize}
  \item DPO/pověřenec: povinný u veřejných orgánů, rozsáhlého pravidelného monitorování nebo rozsáhlého zpracování citlivých údajů.
  \item Práva: informace, přístup, oprava, výmaz, omezení, přenositelnost, námitka, nebýt předmětem čistě automatizovaného rozhodování.
  \item Porušení zabezpečení: oznámit ÚOOÚ do 72 hodin, pokud hrozí riziko; subjektům při vysokém riziku.
  \item ÚOOÚ: dozor, metodika, kontroly, pokuty, řešení stížností.
\end{itemize}

\must{Nikdy nezačínej „dáme souhlas“. Začni účelem, právním titulem, zásadami, minimalizací, rizikem/DPIA a dokumentací.}

\newpage

\otazka{23. Elektronický podpis a elektronická pečeť}{Právní úprava a druhy. Datové schránky, právní úprava a praxe používání.}

\blok{Právní úprava}
\begin{itemize}
  \item eIDAS a eIDAS 2.0: EU rámec pro elektronickou identifikaci a služby vytvářející důvěru.
  \item Zákon o službách vytvářejících důvěru pro elektronické transakce.
  \item Zákon o elektronické identifikaci.
  \item Datové schránky upravuje zákon o elektronických úkonech a autorizované konverzi dokumentů.
  \item DIA: koordinace digitálních služeb státu, datových schránek a příprava evropské peněženky digitální identity v ČR.
\end{itemize}

\blok{Elektronický podpis}
\begin{itemize}
  \item Prostý elektronický podpis: jméno v e-mailu, kliknutí, sken podpisu; právní účinky má, důkazně je slabší.
  \item Zaručený elektronický podpis: kryptografická vazba na podepisujícího, detekce změn, kontrola podepisujícím.
  \item Zaručený podpis s kvalifikovaným certifikátem: ověřená identita u kvalifikovaného poskytovatele; v ČR součást uznávaného podpisu.
  \item Kvalifikovaný elektronický podpis QES: zaručený podpis + kvalifikovaný certifikát + kvalifikovaný prostředek; účinky vlastnoručního podpisu v celé EU.
  \item Uznávaný elektronický podpis v ČR není totéž co QES: každý QES je uznávaný, ne každý uznávaný je QES.
\end{itemize}

\blok{Certifikát prakticky}
\begin{itemize}
  \item Self-signed certifikát si můžeš vytvořit sám, ale není kvalifikovaný.
  \item Kvalifikovaný certifikát vydává kvalifikovaný poskytovatel na důvěryhodném seznamu.
  \item Poskytovatel ověřuje identitu, certifikát váže veřejný klíč na osobu a obsahuje údaje pro ověření a zneplatnění.
  \item Pro QES nestačí certifikát; podpis musí vzniknout kvalifikovaným prostředkem, např. certifikovaný token/karta nebo vzdálená kvalifikovaná služba.
\end{itemize}

\blok{Elektronická pečeť}
\begin{itemize}
  \item Pečeť je pro právnickou osobu, podpis pro fyzickou osobu.
  \item Dokazuje původ a integritu dokumentu, ne vůli konkrétního člověka.
  \item Vhodná pro automatizované dokumenty úřadu nebo firmy.
\end{itemize}

\blok{Časové razítko a validační služby}
\begin{itemize}
  \item Časové razítko dokazuje, že dokument existoval v určitém čase a od té doby se nezměnil.
  \item Důležité pro dlouhodobé ověřování podpisu: certifikát mohl později expirovat nebo být zneplatněn.
  \item Kvalifikované služby vytvářející důvěru poskytují certifikáty, razítka, validaci a uchovávání podpisů/pečetí.
\end{itemize}

\blok{Datové schránky}
\begin{itemize}
  \item Důvěryhodný kanál elektronické komunikace se státem a mezi subjekty.
  \item Doručení má právní účinky; existuje fikce doručení.
  \item Odeslání z datové schránky často nahrazuje potřebu podpisu podání, ale nevytváří podpis uvnitř PDF.
  \item Uchovat datovou zprávu, doručenku a kontext; samotná příloha může ztratit důkazní souvislosti.
  \item Datová schránka identifikuje schránku/subjekt, elektronický podpis identifikuje podepisující osobu a integritu dokumentu.
\end{itemize}

\blok{Autorizovaná konverze}
\begin{itemize}
  \item Převod listina--elektronický dokument nebo elektronický dokument--listina s právními účinky.
  \item Konverze potvrzuje shodu výstupu se vstupem, ne pravdivost obsahu.
\end{itemize}

\must{Závěr: podpis řeší osobu a integritu, pečeť původ právnické osoby, razítko čas, datová schránka doručení a kanál, konverze převod formy.}

\end{document}
